\section{Norms and multiplicities on arbitrary finite free etale fibres}
\label{sec:general-etale-norms}

Let $A\to B$ be a finite free etale homomorphism of commutative rings
and let $x\in\mathcal S(A)$. No trivially valued coefficient field is
assumed. The completed field $H(x)$ and the actual completed tensor
fibre give the norm and multiplicity calculations below.

\begin{lemma}[Scalar and polynomial norms commute with base change]
\label{spectrum:algebra-norm-base-change}
\lean{EnvelopeProperty.normForm,
EnvelopeProperty.eval_map_normForm,
EnvelopeProperty.algebraNorm_baseChange_tmul}
\leanok
\uses{spectrum:base-change-restriction}
For every commutative $A$-algebra $K$ and $b\in B$,
$N_{B/A}(b)$ maps to $N_{K\otimes_A B/K}(1\otimes b)$.
The norm form of $\sum_j X_jb_j$ also commutes with every scalar
extension and polynomial evaluation.
\end{lemma}
\begin{proof}
Choose a finite free basis of $B$ over $A$. The multiplication matrix
for $1\otimes b$ in the scalar-extended basis is obtained by mapping
every coefficient of the original multiplication matrix. Determinants
commute with ring homomorphisms, proving the scalar assertion.
For the polynomial norm form, perform the same matrix calculation
with the multiplication operator of $\sum_j X_j\otimes b_j$ and
then apply polynomial evaluation to every matrix entry. This gives
the actual polynomial base-change and evaluation assertion.
\end{proof}

\begin{lemma}[Scalar field norm and the unique extending seminorm]
\label{spectrum:scalar-field-norm}
\lean{EnvelopeProperty.valuation_algebraNorm_eq_pow,
EnvelopeProperty.valuation_algebraNorm_eq_prod_of_algEquiv}
\leanok
\uses{spectrum:complete-norm-extension}
If $L/H$ is a finite separable extension of a complete ultrametric
normed field and $w$ is its unique extending seminorm, then
$|N_{L/H}(b)|=w(b)^{[L:H]}$. This statement requires no infinitude
hypothesis on $H$.
\end{lemma}
\begin{proof}
Extend the norm of $H$ to its algebraic closure. The field norm is
the product of the images of $b$ under all $H$-embeddings of $L$
into this algebraic closure. Each embedding pulls the extended norm
back to $w$, by genuine uniqueness of complete-field norm extensions.
Multiplicativity therefore makes the product equal to $w(b)$ raised
to the number of embeddings. Separability identifies that number
with $[L:H]$. Norms of finite products are products of norms,
so the same computation applies to a product of finite separable fields.
\end{proof}

\begin{theorem}[Etale fibre multiplicities]
\label{spectrum:etale-fibre-multiplicities}
\lean{EnvelopeProperty.Berkovich.Spectrum.fibreMult,
EnvelopeProperty.Berkovich.Spectrum.one_le_fibreMult,
EnvelopeProperty.Berkovich.Spectrum.fibreEquivPi,
EnvelopeProperty.Berkovich.Spectrum.sum_fibreMult}
\leanok
\uses{spectrum:integral-completed-fibres,spectrum:finite-completed-fibres}
The finite etale tensor algebra $E_x=H(x)\otimes_A B$ is a reduced
Artinian ring and is the product of its maximal quotient fields.
For a point $y$ of the actual fibre, define
$m_y=[E_x/\ker(E_x\to H(y)):H(x)]$.
These residue degrees are positive, and
$\sum_y m_y=\operatorname{rank}_A B$.
\end{theorem}
\begin{proof}
Etale base change makes $E_x$ reduced; module finiteness makes it
Artinian. The reduced Artinian decomposition identifies it with
the product of its maximal quotients, which are finite separable
extensions of $H(x)$. Every such quotient is a nonzero vector space,
giving positivity. Vector-space dimension of the product is the
sum of the quotient dimensions. Finite free base change identifies
the dimension of $E_x$ with the rank of $B$ over $A$; the genuine
maximal tensor support classification reindexes the sum by the
actual seminorm fibre points.
\end{proof}

\begin{theorem}[Scalar norm identity at an arbitrary point]
\label{spectrum:scalar-etale-norm}
\lean{EnvelopeProperty.Berkovich.Spectrum.val_algebraNorm}
\leanok
\uses{spectrum:algebra-norm-base-change,spectrum:scalar-field-norm,
spectrum:etale-fibre-multiplicities}
For every $b\in B$,
\[
 |N_{B/A}(b)|_x=\prod_{y\text{ over }x}|b|_y^{m_y}.
\]
There is no infinitude hypothesis on the completed residue field.
\end{theorem}
\begin{proof}
Base change replaces the left side by the norm of $1\otimes b$
in $E_x$. Its reduced Artinian decomposition turns that norm into
the product of the field norms of its components. The actual
completed-field tensor maps give extending valuations on these
quotient fields, taking value $|b|_y$ on the corresponding component.
Apply the scalar field norm formula and then the genuine fibre
bijection to reindex the finite product.
\end{proof}

\begin{theorem}[Weighted polynomial norm identity]
\label{spectrum:weighted-etale-norm}
\lean{EnvelopeProperty.Berkovich.Spectrum.wgaussValuation_map_normForm,
EnvelopeProperty.Berkovich.Spectrum.wgaussValuation_normForm}
\leanok
\uses{spectrum:etale-fibre-multiplicities,spectrum:algebra-norm-base-change}
Assume explicitly that $H(x)$ is infinite. Let
$P=N(\sum_j X_jb_j)=\sum_\alpha c_\alpha X^\alpha$ and $r_j\ge0$.
Then
\[
 \max_\alpha |c_\alpha|_x r^\alpha
 =\prod_{y\text{ over }x}\left(\max_j |b_j|_y r_j\right)^{m_y}.
\]
\end{theorem}
\begin{proof}
After coefficient base change to $H(x)$, the norm form evaluates
to the norm form of $E_x$ at every $H(x)$-point. Decompose $E_x$
into its finite separable field factors. Over an algebraic closure,
the field norm form equals the product of its conjugate linear forms:
their evaluations agree, and infinitude of $H(x)$ proves equality
of the entire polynomials. Weighted Gauss multiplicativity and
complete-field norm uniqueness give each field-factor contribution
as the maximum linear value raised to its residue degree. Reindex
the product by the actual fibre. Coefficient seminorm preservation
under $A\to H(x)$ returns the left side to the original coefficients.
\end{proof}

This evaluation-characterized producer retains the explicit infinitude
hypothesis needed by its interpolation proof. The actual determinant
norm-form producer in Section~\ref{subsec:general-weighted-norms} removes
that hypothesis through genuine polynomial base change and tensor splitting. Nonetale algebras require local-length multiplicity factors,
and no residue-degree-only norm formula is asserted for them.
